% =====================================================================
%  preamble.tex  —  导言区（SY03-22 Dan Xie《General Argyres-Douglas Theory》中文译注本）
%  由 Lead 维护；译者只读。缺宏时在文末「追加区」append \newcommand，
%  并在回报里列出新增宏名（不得改动上面的已有定义）。
% =====================================================================
\usepackage[a4paper,margin=2.6cm,top=3.0cm,bottom=2.8cm]{geometry}
\usepackage{amsmath,amssymb,amsthm,mathtools}
\usepackage{mathrsfs}
\usepackage{graphicx}
\usepackage{tikz}
\usetikzlibrary{calc,positioning,decorations.pathmorphing,decorations.markings,patterns,arrows}
\usepackage{array,booktabs,longtable,multirow}
\usepackage{caption}
\usepackage{enumitem}
\usepackage{xcolor}
\usepackage{fancyhdr}

% ---------- 字体（Windows，显式路径） ----------
\setCJKmainfont[Path=C:/Windows/Fonts/,BoldFont=simhei.ttf,ItalicFont=simkai.ttf]{simsun.ttc}
\setCJKsansfont[Path=C:/Windows/Fonts/]{simhei.ttf}
\setCJKmonofont[Path=C:/Windows/Fonts/]{simfang.ttf}
\setCJKfamilyfont{zhkai}[Path=C:/Windows/Fonts/]{simkai.ttf}
\setCJKfamilyfont{zhhei}[Path=C:/Windows/Fonts/]{simhei.ttf}
\setCJKfamilyfont{zhfang}[Path=C:/Windows/Fonts/]{simfang.ttf}
\providecommand{\kaishu}{\CJKfamily{zhkai}}
\providecommand{\heiti}{\CJKfamily{zhhei}}
\providecommand{\fangsong}{\CJKfamily{zhfang}}

% ---------- 数学宏（块文件里只许用这些，禁止自造） ----------
\newcommand{\sN}{\mathcal{N}}          % 超对称数 \sN=2
\newcommand{\sA}{\mathcal{A}}          % 平坦联络 𝒜 = A + iφ
\newcommand{\sM}{\mathcal{M}}          % 模空间
\newcommand{\sO}{\mathcal{O}}          % 算子
\newcommand{\sD}{\mathcal{D}}          % 判别式
\DeclareMathOperator{\Tr}{Tr}
\DeclareMathOperator{\diag}{diag}
\DeclareMathOperator{\SU}{SU}
\DeclareMathOperator{\SO}{SO}
\DeclareMathOperator{\U}{U}
\DeclareMathOperator{\Sl}{Sl}
\DeclareMathOperator{\SlN}{Sl}
\newcommand{\R}{\mathbb{R}}
\newcommand{\C}{\mathbb{C}}
\newcommand{\Z}{\mathbb{Z}}
\newcommand{\Q}{\mathbb{Q}}
\newcommand{\CP}{\mathbb{CP}}
\newcommand{\eps}{\varepsilon}
\newcommand{\dd}{\mathrm{d}}
\newcommand{\ee}{\mathrm{e}}
\newcommand{\ii}{\mathrm{i}}
\newcommand{\half}{\tfrac{1}{2}}
\newcommand{\third}{\tfrac{1}{3}}
\newcommand{\Ncoupling}{N_{\text{coupling}}}
% 表格里常用的列类型
\newcolumntype{C}[1]{>{\centering\arraybackslash}p{#1}}

% ---------- 定理类环境（编号一律手工给：\begin{theorem}[3.1]） ----------
\theoremstyle{definition}
\newtheorem*{theorem}{定理}
\newtheorem*{proposition}{命题}
\newtheorem*{lemma}{引理}
\newtheorem*{corollary}{推论}
\newtheorem*{definition}{定义}
\newtheorem*{example}{例}
\newtheorem*{remark}{注}

% ---------- 译注 / 路线注 ----------
% 译注：解释性补充，必须带 [来源：…]（A/B/I 级）或 [译者按（无外部依据）…]
\newcommand{\zhnote}[1]{\par\smallskip\noindent
  {\small\kaishu\color{black!85}【译注】#1}\par\smallskip}
% 路线注：说明作者的论证路线（为什么先证这个 / 哪条路被放弃），必须带原书页码 p.NN
\newcommand{\routenote}[1]{\par\smallskip\noindent
  {\small\kaishu\color{black!70}【路线】#1}\par\smallskip}
% 原文笔误提示（比译注轻）
\newcommand{\typo}[1]{\par\smallskip\noindent
  {\small\kaishu\color{black!85}【原文如此】#1}\par\smallskip}

% ---------- 版式 ----------
\setlength{\parindent}{2em}
\linespread{1.18}
\setlength{\emergencystretch}{2.5em}
\setlist{nosep,leftmargin=2.2em}
\renewcommand{\arraystretch}{1.15}
\captionsetup{font=small,labelsep=quad}

\pagestyle{fancy}
\fancyhf{}
\fancyhead[L]{\small\kaishu 广义 Argyres--Douglas 理论}
\fancyhead[R]{\small\thepage}
\renewcommand{\headrulewidth}{0.4pt}
\fancypagestyle{plain}{\fancyhf{}\fancyfoot[C]{\small\thepage}\renewcommand{\headrulewidth}{0pt}}

\ctexset{
  section = {
    format = \Large\bfseries\heiti,
    beforeskip = 2.2ex plus 1ex minus .2ex,
    afterskip  = 1.4ex plus .2ex,
    name = {},
    number = \arabic{section},
  },
  subsection = {
    format = \large\bfseries\heiti,
    beforeskip = 1.8ex plus .6ex minus .2ex,
    afterskip  = 1.0ex plus .2ex,
  },
  subsubsection = {
    format = \normalsize\bfseries\heiti,
    beforeskip = 1.4ex plus .5ex minus .2ex,
    afterskip  = 0.8ex plus .2ex,
  },
}

\usepackage[hidelinks,bookmarksnumbered=true,bookmarksopen=true,
            pdftitle={广义 Argyres-Douglas 理论（中文译注本）},
            pdfauthor={Dan Xie}]{hyperref}

% ===================== 追加区（译者可 append \newcommand） =====================
